\documentclass[11pt,a4paper]{article}

% ── Packages ───────────────────────────────────────────────────
\usepackage[utf8]{inputenc}
\usepackage[T1]{fontenc}
\usepackage{lmodern}
\usepackage[margin=2.5cm]{geometry}
\usepackage{graphicx}
\usepackage{xcolor}
\usepackage{hyperref}
\usepackage{listings}
\usepackage{booktabs}
\usepackage{enumitem}
\usepackage{titlesec}
\usepackage{fancyhdr}
\usepackage{tcolorbox}
\usepackage{parskip}

% ── Colours ────────────────────────────────────────────────────
\definecolor{fridaybrand}{HTML}{1E5A96}
\definecolor{codebg}{HTML}{F5F2EB}
\definecolor{codefg}{HTML}{23211E}
\definecolor{rulegreen}{HTML}{2E7D32}
\definecolor{rulered}{HTML}{C62828}
\definecolor{sectionblue}{HTML}{1565C0}

% ── Hyperlinks ─────────────────────────────────────────────────
\hypersetup{
  colorlinks  = true,
  linkcolor   = fridaybrand,
  urlcolor    = fridaybrand,
  citecolor   = fridaybrand,
}

% ── Code listings ──────────────────────────────────────────────
\lstset{
  backgroundcolor = \color{codebg},
  basicstyle      = \ttfamily\small\color{codefg},
  breaklines      = true,
  frame           = single,
  rulecolor       = \color{fridaybrand!30},
  numbers         = left,
  numberstyle     = \tiny\color{gray},
  tabsize         = 4,
  showstringspaces= false,
  captionpos      = b,
}

% ── Section styling ────────────────────────────────────────────
\titleformat{\section}
  {\Large\bfseries\color{sectionblue}}{\thesection}{1em}{}
\titleformat{\subsection}
  {\large\bfseries\color{fridaybrand}}{\thesubsection}{1em}{}
\titleformat{\subsubsection}
  {\normalsize\bfseries\color{fridaybrand!80}}{\thesubsubsection}{1em}{}

% ── Header / Footer ───────────────────────────────────────────
\pagestyle{fancy}
\fancyhf{}
\fancyhead[L]{\textcolor{fridaybrand}{\textbf{Project Friday}}}
\fancyhead[R]{\textcolor{gray}{Architecture Document}}
\fancyfoot[C]{\thepage}
\renewcommand{\headrulewidth}{0.4pt}

% ── Callout boxes ──────────────────────────────────────────────
\newtcolorbox{notebox}{
  colback=blue!5, colframe=fridaybrand, title=Note,
  fonttitle=\bfseries, arc=2mm
}
\newtcolorbox{warningbox}{
  colback=red!5, colframe=rulered, title=Warning,
  fonttitle=\bfseries, arc=2mm
}
\newtcolorbox{tipbox}{
  colback=green!5, colframe=rulegreen, title=Tip,
  fonttitle=\bfseries, arc=2mm
}

% ══════════════════════════════════════════════════════════════
\begin{document}

% ── Title ──────────────────────────────────────────────────────
\begin{titlepage}
\centering
\vspace*{3cm}
{\Huge\bfseries\color{fridaybrand} Project Friday\par}
\vspace{0.6cm}
{\Large Architecture \& Design Document\par}
\vspace{1.5cm}
{\large A self-improving, tool-using AI agent\\with a ReAct brain, sandboxed OS hands,\\and a growing skill library.\par}
\vspace{2cm}
{\normalsize Version 0.2.0\par}
{\normalsize \today\par}
\vspace{1cm}
{\small Stack: Python 3.11+ $\cdot$ FastAPI $\cdot$ LangGraph $\cdot$ LangChain $\cdot$ Next.js 15 $\cdot$ React 19 $\cdot$ TypeScript\par}
\vfill
\end{titlepage}

\tableofcontents
\newpage

% ══════════════════════════════════════════════════════════════
\section{Vision \& Scope}

Friday is a personal agentic AI assistant that reasons over a set of tools to complete tasks---from fetching recent CVEs to scaffolding a Next.js component.  It operates on a \textbf{ReAct loop} (Reason $\to$ Act $\to$ Observe $\to$ Reason again) rather than a linear pipeline, meaning it genuinely decides what to do next instead of following a hardcoded flow.

\subsection{Design Philosophy}

\begin{enumerate}[leftmargin=*]
  \item \textbf{Agent, not chatbot.}  Every response is the result of tool-augmented reasoning.
  \item \textbf{Self-improving.}  The Voyager-inspired skill library lets Friday write, test, and commit reusable scripts and framework agents.
  \item \textbf{Security by default.}  All OS access is sandboxed; shell commands are allow\-listed; external directories require explicit opt-in.
  \item \textbf{Local-first.}  LM~Studio is the default provider---no cloud API keys required to start.
\end{enumerate}

% ══════════════════════════════════════════════════════════════
\section{System Architecture}

\begin{lstlisting}[language={},title={High-level data flow}]
+--------------------------------------------------------+
|                  Next.js Frontend                       |
|  Chat UI --SSE Stream--> /api/friday (route handler)   |
+----------------------------+---------------------------+
                             | HTTP POST
+----------------------------v---------------------------+
|                  FastAPI Backend                        |
|  POST /chat --> LangGraph executor --SSE--> client     |
|  GET  /workspace --> file tree snapshot                 |
|  GET  /skills    --> registered skill index             |
|  GET  /agents    --> registered Skill Agents            |
+----------------------------+---------------------------+
                             |
+----------------------------v---------------------------+
|               LangGraph State Machine                   |
|                                                         |
|   +----------+  tool_call?  +--------------------+     |
|   |  Agent   | --- YES ---> |  Tool Executor     |     |
|   |  Node    | <----------- |  Node (cached)     |     |
|   |  (LLM)  |              +--------------------+     |
|   +----+-----+                                         |
|        | NO (final answer)                              |
|        v                                                |
|      END                                                |
+---------------------------------------------------------+
\end{lstlisting}

\subsection{Component Responsibilities}

\begin{table}[h!]
\centering
\begin{tabular}{@{}lp{9cm}@{}}
\toprule
\textbf{Component} & \textbf{Responsibility} \\
\midrule
\texttt{main.py} & FastAPI application, routes, SSE streaming, CORS \\
\texttt{agent/graph.py} & LangGraph \texttt{StateGraph} definition, cached tool node \\
\texttt{agent/nodes.py} & Agent node (LLM invocation with system prompt), router, retry logic \\
\texttt{agent/state.py} & \texttt{AgentState} TypedDict flowing through every node \\
\texttt{agent/system\_prompt.py} & Friday's identity, ReAct instructions, dynamic context injection \\
\texttt{agent/model.py} & Multi-provider LLM factory (LM~Studio, Groq, Ollama) \\
\texttt{agent/tools/} & All tool implementations (see Section~\ref{sec:tools}) \\
\bottomrule
\end{tabular}
\end{table}

% ══════════════════════════════════════════════════════════════
\section{Agent State}

The \texttt{AgentState} TypedDict flows through every node:

\begin{table}[h!]
\centering
\begin{tabular}{@{}llp{7cm}@{}}
\toprule
\textbf{Field} & \textbf{Type} & \textbf{Purpose} \\
\midrule
\texttt{messages} & \texttt{list} & Conversation history (LangGraph's \texttt{add\_messages} reducer) \\
\texttt{tool\_attempts} & \texttt{int} & Retry counter for the current tool/script \\
\texttt{active\_skill} & \texttt{str|None} & Path of the skill script being executed \\
\texttt{skill\_context} & \texttt{str|None} & Loaded Skill Agent style guide (injected into system prompt) \\
\texttt{target\_directory} & \texttt{str|None} & External project directory (must be in \texttt{ALLOWED\_TARGET\_DIRS}) \\
\texttt{error\_history} & \texttt{list[str]} & Rolling log of error messages for retry decisions \\
\bottomrule
\end{tabular}
\end{table}

% ══════════════════════════════════════════════════════════════
\section{Tools}
\label{sec:tools}

\subsection{File Operations}

\begin{table}[h!]
\centering
\begin{tabular}{@{}lp{9cm}@{}}
\toprule
\textbf{Tool} & \textbf{Description} \\
\midrule
\texttt{list\_files} & Lists files/folders at a given workspace-relative path \\
\texttt{read\_file} & Reads full contents of a workspace file \\
\texttt{write\_to\_file} & Creates or overwrites a workspace file \\
\texttt{append\_to\_file} & Appends content to an existing file (or creates it) \\
\texttt{create\_directory} & Creates a directory (and parents) inside the workspace \\
\texttt{delete\_file} & Deletes a single file (not directories) from the workspace \\
\bottomrule
\end{tabular}
\end{table}

\subsection{Shell Execution}

\begin{table}[h!]
\centering
\begin{tabular}{@{}lp{9cm}@{}}
\toprule
\textbf{Tool} & \textbf{Description} \\
\midrule
\texttt{execute\_bash\_command} & Runs an allowlisted command inside the workspace sandbox \\
\texttt{execute\_in\_directory} & Runs an allowlisted command in an external allowed directory \\
\bottomrule
\end{tabular}
\end{table}

\begin{notebox}
Shell commands are validated against a configurable \textbf{allowlist} (first token must match).  Shell chaining operators (\texttt{\&\&}, \texttt{||}, \texttt{;}, \texttt{|}) are rejected.  Timeout defaults to 120 seconds.
\end{notebox}

\subsection{Web Search}

\texttt{web\_search} queries the web via Tavily (if configured) or DuckDuckGo (fallback), returning the top 3 result summaries.

\subsection{Skill Library}

\begin{table}[h!]
\centering
\begin{tabular}{@{}lp{9cm}@{}}
\toprule
\textbf{Tool} & \textbf{Description} \\
\midrule
\texttt{save\_to\_skill\_library} & Moves a tested script from workspace into the skills directory \\
\texttt{search\_skill\_library} & Finds the closest matching skill by description similarity \\
\bottomrule
\end{tabular}
\end{table}

\subsection{Skill Agents}

\begin{table}[h!]
\centering
\begin{tabular}{@{}lp{9cm}@{}}
\toprule
\textbf{Tool} & \textbf{Description} \\
\midrule
\texttt{create\_skill\_agent} & Creates a new framework-specific Skill Agent with manifest, style guide, and templates directory \\
\texttt{load\_skill\_context} & Loads an agent's manifest and style guide into the LLM context \\
\texttt{list\_skill\_agents} & Lists all registered Skill Agents and their descriptions \\
\texttt{update\_skill\_agent} & Patches an existing agent's rules, description, or scaffold steps \\
\bottomrule
\end{tabular}
\end{table}

% ══════════════════════════════════════════════════════════════
\section{Skill Agent System}

The Skill Agent system is the core differentiator---it allows Friday to \textbf{teach itself} new frameworks and domains at runtime.

\subsection{Directory Structure}

\begin{lstlisting}[language={},title={Skill Agent layout}]
skills/
  index.json              # Script skill registry
  agents/
    nextjs/
      manifest.json       # Metadata, triggers, dos/don'ts
      style_guide.md      # Auto-generated best-practice doc
      templates/          # Optional starter files
    fastapi/
      manifest.json
      style_guide.md
      templates/
\end{lstlisting}

\subsection{Manifest Schema}

Each Skill Agent is defined by a \texttt{manifest.json} with the following fields:

\begin{table}[h!]
\centering
\begin{tabular}{@{}llp{6.5cm}@{}}
\toprule
\textbf{Field} & \textbf{Type} & \textbf{Description} \\
\midrule
\texttt{name} & string & Slugified agent name (e.g.\ \texttt{nextjs}) \\
\texttt{display\_name} & string & Human-readable name (e.g.\ \texttt{Next.js}) \\
\texttt{version} & string & Semantic version \\
\texttt{description} & string & What this agent handles \\
\texttt{trigger\_patterns} & string[] & Keywords that activate this agent \\
\texttt{rules.do} & string[] & Best-practice rules to follow \\
\texttt{rules.dont} & string[] & Anti-patterns to avoid \\
\texttt{scaffold\_steps} & string[] & Shell commands to set up a new project \\
\texttt{created\_at} & ISO 8601 & Creation timestamp \\
\texttt{updated\_at} & ISO 8601 & Last modification timestamp \\
\texttt{auto\_generated} & boolean & Whether the LLM created this agent \\
\bottomrule
\end{tabular}
\end{table}

\subsection{Lifecycle}

\begin{enumerate}[leftmargin=*]
  \item \textbf{Detection:} User asks to work with a framework (e.g.\ ``Create a Next.js project'').
  \item \textbf{Lookup:} Agent calls \texttt{search\_skill\_library("next.js")}.
  \item \textbf{Creation (first time):} No match found $\to$ agent calls \texttt{create\_skill\_agent(...)} with framework-specific dos, don'ts, scaffold steps, and trigger patterns.
  \item \textbf{Loading:} Agent calls \texttt{load\_skill\_context("nextjs")} $\to$ style guide is injected into the system prompt.
  \item \textbf{Execution:} Agent follows the loaded rules strictly while scaffolding the project.
  \item \textbf{Evolution:} If the agent learns new patterns or encounters errors, it calls \texttt{update\_skill\_agent(...)} to refine the rules.
\end{enumerate}

\begin{tipbox}
On subsequent requests involving the same framework, the agent skips creation and immediately loads the existing agent's context---making it faster and more consistent over time.
\end{tipbox}

% ══════════════════════════════════════════════════════════════
\section{Security Model}

\subsection{Path Sandboxing}

All file operations go through \texttt{safe\_path()}, which enforces:

\begin{itemize}[leftmargin=*]
  \item \textbf{No absolute paths} inside the workspace.
  \item \textbf{No traversal} via \texttt{../../} or similar.
  \item \textbf{Symlink detection:} if resolving a path follows a symlink, the target must still be inside the workspace.
  \item \textbf{Blocked segments:} paths containing \texttt{.env}, \texttt{.ssh}, \texttt{.git}, \texttt{.gnupg}, or \texttt{.aws} are rejected.
\end{itemize}

For operations outside the workspace, \texttt{safe\_target\_path()} validates against \texttt{ALLOWED\_TARGET\_DIRS}.

\subsection{Command Allowlist}

The blocklist-based approach was replaced with an \textbf{allowlist}:

\begin{itemize}[leftmargin=*]
  \item Only the first token of a command is checked against a configurable list of permitted prefixes.
  \item Shell chaining operators (\texttt{\&\&}, \texttt{||}, \texttt{;}, \texttt{|}, backticks, \texttt{\$()}) are rejected.
  \item The allowlist is stored in the \texttt{COMMAND\_ALLOWLIST} environment variable.
\end{itemize}

\begin{warningbox}
The command is still executed with \texttt{shell=True} for compatibility with complex commands.  For maximum security in production, consider switching to container-based isolation (Docker, Firejail, or Windows Sandbox).
\end{warningbox}

\subsection{Skill Script Execution}

Committed skill scripts are executed with \texttt{subprocess.run([sys.executable, path, ...args])}---\textbf{no shell injection} is possible because:

\begin{itemize}[leftmargin=*]
  \item Arguments are passed as a list, not interpolated into a string.
  \item \texttt{shell=True} is NOT used for skill execution.
  \item The Python interpreter is resolved via \texttt{sys.executable}.
\end{itemize}

\subsection{Retry \& Recursion Limits}

\begin{table}[h!]
\centering
\begin{tabular}{@{}llp{7cm}@{}}
\toprule
\textbf{Limit} & \textbf{Default} & \textbf{Purpose} \\
\midrule
\texttt{MAX\_TOOL\_ATTEMPTS} & 3 & Max retries for a failing tool before bailing \\
\texttt{RECURSION\_LIMIT} & 25 & Hard cap on LangGraph cycles per request \\
\texttt{MAX\_TOKENS} & 2048 & Max tokens per LLM response \\
\texttt{COMMAND\_TIMEOUT} & 120s & Max runtime for any subprocess \\
\bottomrule
\end{tabular}
\end{table}

% ══════════════════════════════════════════════════════════════
\section{System Prompt}

The agent receives a comprehensive system prompt on every invocation, assembled by \texttt{build\_system\_prompt()}:

\begin{enumerate}[leftmargin=*]
  \item \textbf{Static core:} Identity, ReAct instructions, tool usage rules, skill agent workflow, self-improvement protocol, and safety constraints.
  \item \textbf{Dynamic skill context:} When a Skill Agent is loaded, its style guide (dos, don'ts, scaffold steps) is appended as an additional section.
  \item \textbf{Extra instructions:} Per-task overrides can be injected programmatically.
\end{enumerate}

This ensures the LLM always knows:
\begin{itemize}[leftmargin=*]
  \item It is called ``Friday'' and operates as an agent, not a chatbot.
  \item Which tools are available and when to use each.
  \item How to find, create, and follow Skill Agents.
  \item Its safety boundaries.
\end{itemize}

% ══════════════════════════════════════════════════════════════
\section{LLM Provider Configuration}

\begin{table}[h!]
\centering
\begin{tabular}{@{}llp{6cm}@{}}
\toprule
\textbf{Provider} & \textbf{Model} & \textbf{Use Case} \\
\midrule
LM~Studio (local) & auto-detected & Default---fast, free, offline \\
Groq & \texttt{llama-3.3-70b-versatile} & Speed-optimized cloud fallback \\
Ollama & \texttt{qwen2.5-coder:7b} & Offline code generation \\
\bottomrule
\end{tabular}
\end{table}

The provider is selected via \texttt{MODEL\_PROVIDER} in \texttt{.env}.  LM~Studio uses auto-discovery (\texttt{/v1/models} endpoint) to find loaded models.

% ══════════════════════════════════════════════════════════════
\section{Frontend Architecture}

\subsection{Component Tree}

\begin{lstlisting}[language={},title={React component hierarchy}]
RootLayout (layout.tsx)
  HomePage (page.tsx)
    FridayChat
      MessageBubble (per event)
      ThinkingIndicator (during streaming)
      WorkspacePanel | SkillsPanel (sidebar tabs)
\end{lstlisting}

\subsection{SSE Stream Hook}

\texttt{useFridayStream()} manages the full lifecycle:

\begin{enumerate}[leftmargin=*]
  \item Creates a conversation ID (persists across messages in one session).
  \item Sends the query via \texttt{POST /api/friday}.
  \item Reads the SSE stream via \texttt{ReadableStream} API.
  \item Deduplicates events to prevent repeated rendering.
  \item Supports abort (Stop button) and reset (New Chat).
\end{enumerate}

\subsection{Event Types}

\begin{table}[h!]
\centering
\begin{tabular}{@{}llp{6cm}@{}}
\toprule
\textbf{Type} & \textbf{Source} & \textbf{UI Rendering} \\
\midrule
\texttt{user} & Frontend & Right-aligned blue bubble \\
\texttt{thought} & Agent reasoning & Dimmed dashed-border card with ``Reasoning step'' badge \\
\texttt{tool\_call} & Agent $\to$ tool & Dashed card with ``Tool call'' badge and JSON args \\
\texttt{tool\_result} & Tool $\to$ agent & Collapsible \texttt{<details>} with ``Tool result'' badge \\
\texttt{final} & Agent answer & Full left-aligned message bubble \\
\bottomrule
\end{tabular}
\end{table}

% ══════════════════════════════════════════════════════════════
\section{Project Structure}

\begin{lstlisting}[language={},title={Full monorepo layout}]
Project-Friday/
  README.md
  Project-info.txt
  run-friday.ps1              # One-command launcher
  docs/
    friday_architecture.tex   # This document
    friday_api_reference.tex  # API reference
  friday-backend/
    main.py                   # FastAPI app + routes
    requirements.txt
    .env / .env.example
    agent/
      __init__.py
      graph.py                # LangGraph state machine
      nodes.py                # Agent node, router, retry
      state.py                # AgentState TypedDict
      system_prompt.py        # System prompt builder
      model.py                # LLM provider factory
      tools/
        __init__.py           # Tool registry
        os_tools.py           # File + shell tools
        web_search.py         # Tavily / DuckDuckGo
        skill_library.py      # Save / search skills
        skill_agent.py        # Skill Agent CRUD
    workspace/                # Sandboxed file system
    skills/
      index.json              # Script skill registry
      agents/                 # Skill Agent definitions
  friday-frontend/
    app/
      page.tsx                # Root page
      layout.tsx              # Fonts + metadata
      globals.css             # Design tokens
      api/friday/
        route.ts              # Proxy: chat + workspace
        skills/route.ts       # Proxy: skills
        agents/route.ts       # Proxy: agents
    components/
      FridayChat.tsx          # Main chat container
      MessageBubble.tsx       # Event renderer
      ThinkingIndicator.tsx   # Streaming dots
      WorkspacePanel.tsx      # File tree sidebar
      SkillsPanel.tsx         # Skills + agents sidebar
    hooks/
      useFridayStream.ts      # SSE consumer hook
    lib/
      api.ts                  # Typed API client
\end{lstlisting}

% ══════════════════════════════════════════════════════════════
\section{Future Enhancements}

\begin{enumerate}[leftmargin=*]
  \item \textbf{Semantic skill search:} Replace \texttt{SequenceMatcher} with embedding-based lookup (ChromaDB or LanceDB).
  \item \textbf{Persistent memory:} Replace \texttt{MemorySaver} with SQLite or file-based checkpointer for conversation persistence across restarts.
  \item \textbf{MCP integration:} Connect to Model Context Protocol servers for external tool providers.
  \item \textbf{Multi-language skills:} Support non-Python skill scripts (Node.js, Bash, Go).
  \item \textbf{Container isolation:} Run tool commands in Docker containers for production-grade sandboxing.
  \item \textbf{Authentication:} Add API key or JWT-based auth for production deployments.
\end{enumerate}

\end{document}
